% ==============================================================================
% Chapitre 7 : Annexes
% Ports et protocoles, commandes maintenance, glossaire, références
% ==============================================================================

\chapter{Annexes}

\section{Référence des ports et protocoles}

Le tableau~\ref{tab:ports-protocoles} récapitule l'ensemble des ports et
protocoles utilisés par les différents composants de l'infrastructure
LockBits. Il distingue trois réseaux logiques~: le réseau \texttt{frontend}
(services exposés), le réseau \texttt{backend} (services internes) et le
réseau \texttt{monitoring} (supervision).

\begin{table}[htbp]
  \centering
  \caption{Référence des ports et protocoles}
  \label{tab:ports-protocoles}
  \begin{tabular}{llllp{5cm}}
    \toprule
    \textbf{Service} & \textbf{Port} & \textbf{Protocole} & \textbf{Réseau} & \textbf{Description} \\
    \midrule
    EDR Server       & 8001 & HTTP/TCP   & frontend   & API REST FastAPI \\
    Site Web         & 8080 & HTTP/TCP   & frontend   & Site PHP/Apache \\
    MySQL            & 3306 & MySQL/TCP  & backend    & Base de données du site \\
    Prometheus       & 9090 & HTTP/TCP   & monitoring & Métriques système \\
    Grafana          & 3000 & HTTP/TCP   & monitoring & Dashboards de visualisation \\
    Alloy            & --   & gRPC/HTTP  & monitoring & Collecte et acheminement des logs \\
    VirusTotal       & 443  & HTTPS      & externe     & API d'analyse de fichiers \\
    Claude API       & 443  & HTTPS      & externe     & API du chatbot intelligent \\
    GLPI             & 443  & HTTPS      & externe     & API de gestion de tickets \\
    GHCR             & 443  & HTTPS      & externe     & Registry Docker GitHub \\
    \bottomrule
  \end{tabular}
\end{table}

Les services exposés sur le réseau \texttt{frontend} sont accessibles depuis
l'extérieur via les ports configurés dans le fichier \texttt{.env}. Les
services internes (\texttt{backend}) ne sont accessibles qu'entre conteneurs
Docker, ce qui limite la surface d'attaque. Le réseau \texttt{monitoring}
reste isolé du trafic applicatif.

\section{Commandes de maintenance courantes}

Cette section regroupe les commandes essentielles pour l'administration
quotidienne de l'infrastructure LockBits.

\subsection{Administration Docker}

La stack LockBits est déployée via Docker Compose. Les commandes suivantes
permettent de gérer l'état des services~:

\begin{lstlisting}[language=dockerbash, caption=Voir l'état des services]
docker compose -f docker-compose.prod.yml ps
\end{lstlisting}

\begin{lstlisting}[language=dockerbash, caption=Consulter les logs en temps réel]
docker compose -f docker-compose.prod.yml logs -f
\end{lstlisting}

\begin{lstlisting}[language=dockerbash, caption=Redémarrer un service spécifique]
docker compose -f docker-compose.prod.yml restart edr-server
\end{lstlisting}

\begin{lstlisting}[language=dockerbash, caption=Mettre à jour tous les services]
docker compose -f docker-compose.prod.yml pull
docker compose -f docker-compose.prod.yml up -d
\end{lstlisting}

\begin{lstlisting}[language=dockerbash, caption=Sauvegarder la base MySQL]
docker exec lockbits_db mysqldump -u lockbits -p lockbits_client > backup.sql
\end{lstlisting}

\subsection{Installation et mise à jour (one-liner)}

Le script \texttt{install.sh} centralisé le déploiement et les mises à jour
de la stack. Il est exécutable en une seule commande~:

\begin{lstlisting}[language=dockerbash, caption=Installation complète de la stack]
curl -fsSL https://raw.githubusercontent.com/Lockbits-ESGI/main/main/install.sh | bash
\end{lstlisting}

\begin{lstlisting}[language=dockerbash, caption=Mise à jour de la stack]
curl -fsSL https://raw.githubusercontent.com/Lockbits-ESGI/main/main/install.sh | bash -s -- --update
\end{lstlisting}

\begin{lstlisting}[language=dockerbash, caption=Activer la mise à jour automatique (cron)]
curl -fsSL https://raw.githubusercontent.com/Lockbits-ESGI/main/main/install.sh | bash -s -- --cron
\end{lstlisting}

\subsection{Gestion des sous-modules Git}

Le dépôt principal LockBits orchestre ses composants via des sous-modules
Git. Les commandes suivantes permettent de les maintenir à jour~:

\begin{lstlisting}[language=dockerbash, caption=Cloner le projet avec les sous-modules]
git clone --recurse-submodules git@github.com:Lockbits-ESGI/main.git
\end{lstlisting}

\begin{lstlisting}[language=dockerbash, caption=Mettre à jour les sous-modules et publier]
git submodule update --remote
git add -u
git commit -m "update submodules"
git push
\end{lstlisting}

\begin{lstlisting}[language=dockerbash, caption=Mise à jour via le script dédié]
bash install.sh --update-sources
\end{lstlisting}

\subsection{Agent EDR}

L'agent de sécurité LockBits s'installe sur les machines Linux à protéger~:

\begin{lstlisting}[language=dockerbash, caption=Installation en mode binaire]
curl -fsSL https://raw.githubusercontent.com/Lockbits-ESGI/main/main/install-agent.sh | \
  SERVER_URL=https://edr.lockbits.pro AUTH_TOKEN=token bash
\end{lstlisting}

\begin{lstlisting}[language=dockerbash, caption=Installation en mode Docker]
curl -fsSL https://raw.githubusercontent.com/Lockbits-ESGI/main/main/install-agent.sh | \
  SERVER_URL=https://edr.lockbits.pro AUTH_TOKEN=token bash -s -- --mode docker
\end{lstlisting}

\begin{lstlisting}[language=dockerbash, caption=Désinstallation de l'agent]
sudo /opt/lockbits-agent/lockbits-agent --uninstall
\end{lstlisting}

\section{Glossaire technique}

Le glossaire ci-dessous définit les principaux termes techniques utilisés
dans ce dossier d'architecture.

\begin{longtable}{lp{10cm}}
  \caption{Glossaire technique}
  \label{tab:glossaire} \\
  \toprule
  \textbf{Terme} & \textbf{Définition} \\
  \midrule
  \endfirsthead

  \caption{Glossaire technique (suite)} \\
  \toprule
  \textbf{Terme} & \textbf{Définition} \\
  \midrule
  \endhead

  \midrule
  \multicolumn{2}{r}{\small\textit{Suite sur la page suivante}} \\
  \endfoot

  \bottomrule
  \endlastfoot

  \textbf{EDR}
  & Endpoint Détection \& Response~--~solution de détection et réponse sur
  les terminaux. \\

  \textbf{FIM}
  & File Integrity Monitoring~--~surveillance de l'intégrité des fichiers
  du système. \\

  \textbf{SIEM}
  & Security Information and Event Management~--~gestion des événements et
  informations de sécurité. \\

  \textbf{IaC}
  & Infrastructure as Code~--~gestion et provisionnement de
  l'infrastructure via du code versionné. \\

  \textbf{GHCR}
  & GitHub Container Registry~--~registre de conteneurs intégré à GitHub. \\

  \textbf{CI/CD}
  & Continuous Integration / Continuous Deployment~--~intégration et
  déploiement continus. \\

  \textbf{Cloudflare Zero Trust}
  & Plateforme de sécurité réseau cloud qui fournit un accès sécurisé
  via tunnel cloudflared. \\

  \textbf{MVP}
  & Minimum Viable Product~--~produit minimum viable, première version
  fonctionnelle d'un logiciel. \\

  \textbf{HA}
  & Haute Disponibilité~--~capacité d'un système à rester opérationnel
  malgré des défaillances. \\

  \textbf{SLA}
  & Service Level Agreement~--~accord définissant le niveau de service
  attendu entre un fournisseur et un client. \\

  \textbf{SOC}
  & Security Operations Center~--~centre opérationnel de sécurité
  supervisant les infrastructures. \\

  \textbf{OAuth2}
  & Protocole d'authentification déléguée permettant l'accès sécurisé à
  des ressources sans partage de mot de passe. \\

  \textbf{REST}
  & Representational State Transfer~--~style d'architecture API basé sur
  les verbes HTTP. \\

  \textbf{FastAPI}
  & Framework Python moderne pour la construction d'API asynchrones,
  utilisé pour le serveur EDR. \\

  \textbf{PyInstaller}
  & Outil de compilation permettant de transformer un programme Python en
  binaire exécutable autonome. \\

  \textbf{Trivy}
  & Scanner de vulnérabilités open-source spécialisé dans l'analyse
  d'images conteneurisées. \\

  \textbf{Dokploy}
  & Plateforme de déploiement automatisé pour applications
  conteneurisées, déclenchée par webhook GitHub Actions. \\

  \textbf{Alloy}
  & Agent de collecte et d'acheminement de logs de la stack Grafana
  (anciennement Grafana Agent). \\

  \textbf{STIX}
  & Structured Threat Information eXpression~--~format standardisé
  d'échange de données de menace. \\

\end{longtable}

\section{Références bibliographiques}

Les ressources suivantes ont été consultées pour la conception et la
rédaction de ce dossier d'architecture.

\subsection*{Documentations officielles}

\begin{itemize}
  \item Documentation FastAPI~: \url{https://fastapi.tiangolo.com/}
  \item Documentation Docker Compose~: \url{https://docs.docker.com/compose/}
  \item Documentation GitHub Actions~: \url{https://docs.github.com/en/actions}
  \item Prometheus~: \url{https://prometheus.io/docs/}
  \item Grafana~: \url{https://grafana.com/docs/}
\end{itemize}

\subsection*{API et services externes}

\begin{itemize}
  \item Claude API (Anthropic)~: \url{https://docs.anthropic.com/}
  \item VirusTotal API~: \url{https://developers.virustotal.com/}
  \item GLPI API~: \url{https://glpi-project.org/}
\end{itemize}

\subsection*{Dépôts et organisation}

\begin{itemize}
  \item Organisation GitHub LockBits~: \url{https://github.com/Lockbits-ESGI}
  \item Registre GHCR LockBits~: \url{https://ghcr.io/lockbits-esgi}
\end{itemize}
